\documentclass[12pt]{article}

\setlength{\parindent}{0pt}
\setcounter{secnumdepth}{3}
\usepackage[margin=1in]{geometry}

% packages
\usepackage{amsmath, amsfonts, amsthm, amssymb}
\usepackage{graphicx, float}
\usepackage{mathtools}
\usepackage{titlesec}
\usepackage{interval}
\usepackage{hyperref}
\usepackage{siunitx}
\usepackage{titling}
\usepackage{vwcol}
\usepackage{setspace}
\usepackage{empheq}
\usepackage{cancel}
\usepackage{esdiff}
\usepackage{multicol}
\usepackage{mdframed}
\usepackage{esdiff}
\usepackage{tikzsymbols}
\usepackage{multicol}
\usepackage{tikz}
\usepackage{varwidth}
\usepackage{pgfplots}
\usepackage{fancyhdr}
\usepackage{enumitem}
\usepackage{tcolorbox}

\intervalconfig {
	soft open fences
}

% header
\pagestyle{fancy}
\fancyhf{}
\lhead{Ethan Fan}
\rhead{PCH}
\cfoot{\thepage}

% section format
\titleformat{\section}
  {\bfseries}
  {}
  {0em}
  {}
\titlespacing*{\section}{0pt}{0pt}{0.3em}

\titleformat{\subsection}[runin]
  {\bfseries}
  {}
  {0em}
  {}
\titlespacing*{\subsection}{0pt}{0pt}{0em}

\titleformat{\subsubsection}[runin]
  {\itshape}
  {(\roman{subsubsection})}
  {0em}
  {}

% commands
\newcommand{\alignedintertext}[1]{%
  \noalign{%
    \vskip\belowdisplayshortskip
    \vtop{\hsize=\linewidth#1\par
    \expandafter}%
    \expandafter\prevdepth\the\prevdepth
  }%
}

\newcommand*{\paren}[1]{\ensuremath\left(#1\right)}
\newcommand*{\limit}[2][x]{\ensuremath{\displaystyle\lim_{#1\to#2}}}
\newcommand*{\Limit}[3][x]{\ensuremath{\displaystyle\lim_{#1\to#2}\left[#3\right]}}
\newcommand*{\deriv}[1][x]{\ensuremath{\dfrac{\mathrm{d}}{\mathrm{d}#1}}}
\newcommand*{\Deriv}[2][x]{\ensuremath{\dfrac{\mathrm{d}}{\mathrm{d}#1}\left[#2\right]}}
\newcommand{\dydx}{\frac{dy}{dx}}
\newcommand*{\iinteg}[2][x]{\ensuremath{\displaystyle\int #2\;\mathrm{d}#1}}
\newcommand*{\dinteg}[4][x]{\ensuremath{\displaystyle\int_{#2}^{#3}#4\;\mathrm{d}#1}}
\newcommand*{\abs}[1]{\ensuremath{\left|#1\right|}}
\newcommand*{\eps}{\varepsilon}
\newcommand*{\real}{\ensuremath\mathbb{R}}
\newcommand*{\integer}{\ensuremath\mathbb{Z}}
\newcommand*{\posint}{\ensuremath\mathbb{Z^+}}
\newcommand*{\cbrt}[1]{\ensuremath{\sqrt[3]{#1}}}
\newcommand*{\lheqzero}{\ensuremath{\underset{\text{L'H}}{\overset{\left[\frac00\right]}{=}}}}
\newcommand*{\lheqinfty}{\ensuremath{\underset{\text{L'H}}{\overset{\left[\frac{\infty}{\infty}\right]}{=}}}}
\newcommand{\tor}{\text{ or }}
\newcommand{\floor}[1]{\lfloor #1 \rfloor}

\newcommand{\vem}{\vspace{0.5em}}
\newcommand{\example}[1]{\section{Example #1}}
\newcommand{\problem}[1]{\section{Problem #1}}
\newcommand{\subproblem}[1]{\subsection{(#1)} \,}

\DeclareMathOperator{\dne}{DNE}
\DeclareMathOperator{\sgn}{sgn}

\newcommand{\definition}[1]{\begin{tcolorbox}[colback=red!5!white,colframe=red!75!black,parbox=false] #1 \end{tcolorbox}}

\newcommand{\theorem}[2]{\begin{tcolorbox}[title={#1},colback=blue!5!white,colframe=blue!75!black,parbox=false] #2 \end{tcolorbox}}

\allowdisplaybreaks
% \postdisplaypenalty=100000

% document
\begin{document}

\vspace*{0cm}

% opening
\begin{center}
    {\Large \bfseries Problem Set 64} \\[0.5cm]
    {\large Derivatives 3} \\[0.35cm]
    {\normalsize April 10, 2026}
\end{center}

% problems

\problem{1}
Discuss continuity and differentiability of the first and second derivatives at $x = 0$ in problems 1-4. Justify your answers. \vem 

\subproblem{a} Let
\[
f(x)=\begin{cases}
    x \sin \paren{\dfrac{1}{x}} &x\ne 0\\
    0 &x=0
\end{cases}
\]
\[
\boxed{f(x) $ is continuous$,
    f'(x) \dne }
\]

\subproblem{b}
\[
f(x)=\begin{cases}
    x^2 \sin \paren{\dfrac{1}{x}} &x\ne 0\\
    0 &x=0
\end{cases}
\]
\[
\boxed{f(x) $ is differentiable$,
    f'(x) $ exists but is not continuous$,
    f''(x) \dne}
\]

\subproblem{c}
\[
f(x)=\begin{cases}
    x^3 \sin \paren{\dfrac{1}{x}} &x\ne 0\\
    0 &x=0
\end{cases}
\]
\[
\boxed{f(x) $ is differentiable$,
    f'(x) $ is continuous$,
    f''(x) \dne}
\]

\subproblem{d}
\[
f(x)=\begin{cases}
    x^4 \sin \paren{\dfrac{1}{x}} &x\ne 0\\
    0 &x=0
\end{cases}
\]
\begin{gather*}
    \lim_{h \to 0^-} \frac{f(x+h)-f(x)}{h}=\lim_{h \to 0^-}\frac{h^4 \sin \frac{1}{h}}{h}=\lim_{h \to 0^-}h^3 \sin \frac{1}{h}=0\\
    \lim_{h \to 0^+} \frac{f(x+h)-f(x)}{h}=\lim_{h \to 0^+}\frac{h^4 \sin \frac{1}{h}}{h}=\lim_{h \to 0^+}h^3 \sin \frac{1}{h}=0\\
    f'(0^-)=f'(0^+)=0 \implies f'(0)=0\\
    f(x) $ is differentiable at $ x=0\\
    f'(x)=\begin{cases}
        4x^3 \sin \paren{\dfrac{1}{x}} -x^2 \cos \frac{1}{x} &x\ne 0\\
    0 &x=0
    \end{cases}\\
    \lim_{h \to 0^-} \frac{f'(x+h)-f'(x)}{h}=\lim_{h \to 0^-}\frac{4h^3 \sin \frac{1}{h} -h^2\cos\frac{1}{h}}{h}=\lim_{h \to 0^-} 4h^2 \sin \frac{1}{h}-h \cos \frac{1}{h}=0\\
    \lim_{h \to 0^+} \frac{f'(x+h)-f'(x)}{h}=\lim_{h \to 0^+}\frac{4h^3 \sin \frac{1}{h} -h^2\cos\frac{1}{h}}{h}=\lim_{h \to 0^+} 4h^2 \sin \frac{1}{h}-h \cos \frac{1}{h}=0\\
    f''(0^-)=f''(0^+)=0 \implies f''(0)=0\\
    f'(x)$ is differentiable at $ x=0\\
    f''(x)$ oscillates near 0$\\
\boxed{f(x) $ is differentiable$,
    f'(x) $ is differentiable$,
    f''(x) $ exists but is not continuous$}
\end{gather*}

\problem{2}
Establish a general continuity and differentiability pattern for \vem

\subproblem{a}
\[
f(x)=\begin{cases}
    x^{2n} \sin \paren{ \dfrac{1}{x}} &x\ne 0\\
    0 &x=0
\end{cases}
\]
\begin{minipage}[t]{0.5\linewidth}
    \begin{align*}
        f'(0^-)&=\lim_{h\to0^-} \frac{f(x+h)-f(x)}{h}\\
        &= \lim_{h\to0^-} \frac{h^{2n}\sin \frac{1}{h}}{h}\\
        &=\lim_{h\to0^-} h^{2n-1} \sin \frac{1}{h}\\
        &=0
    \end{align*}
\end{minipage}
\begin{minipage}[t]{0.5 \linewidth}
    \begin{align*}
        f'(0^+)&=\lim_{h\to0^+} \frac{f(x+h)-f(x)}{h}\\
        &= \lim_{h\to0^+} \frac{h^{2n}\sin \frac{1}{h}}{h}\\
        &=\lim_{h\to0^+} h^{2n-1} \sin \frac{1}{h}
        \\&=0
    \end{align*}
\end{minipage}

\begin{gather*}
    f'(0^-)=f'(0^+)=0 \implies f'(0)=0\\
    \boxed{f(x) $ is differentiable at $ x=0 $ for all $ n \ge 1}\\ 
    f'(x)=\begin{cases}
        2nx^{2n-1} \sin \frac{1}{x}-x^{2n-2} \cos \frac{1}{x} &x \ne 0\\
        0 & x=0
    \end{cases}
\end{gather*}
\begin{minipage}[t]{0.5 \linewidth}
    \begin{align*}
        f''(0^-)&=\lim_{h\to0^-} \frac{f'(x+h)-f'(x)}{h}\\
        &= \lim_{h\to0^-}=\frac{2nh^{2n-1}\sin \frac{1}{h}- h^{2n-2}\cos \frac{1}{h}}{h}\\
        &=\lim_{h\to0^-} 2nh^{2n-2} \sin \frac{1}{h}-h^{2n-3}\cos \frac{1}{h}\\
        &=0
    \end{align*}
\end{minipage}
\begin{minipage}[t]{0.5 \linewidth}
    \begin{align*}
        f''(0^+)&=\lim_{h\to0^+} \frac{f'(x+h)-f'(x)}{h}\\
        &= \lim_{h\to0^+}=\frac{2nh^{2n-1}\sin \frac{1}{h}- h^{2n-2}\cos \frac{1}{h}}{h}\\
        &=\lim_{h\to0^+} 2nh^{2n-2} \sin \frac{1}{h}-h^{2n-3}\cos \frac{1}{h}\\
        &=0
    \end{align*}
\end{minipage}

\begin{gather*}
    f''(0^-)=f''(0^+)=0 \implies f''(0)=0\\
    \boxed{f'(x) $ is differentiable at $ x=0 $ for all $ n \ge 2}
\end{gather*}
As the number of times for which $f(x)$ is differentiable increases as $n$ increases, we observe that:
\[
\boxed{f^{(n)}(x) $ exists but is not continuous$}
\]

\subproblem{b}
\[
f(x)=\begin{cases}
    x^{2n+1} \sin \paren{ \dfrac{1}{x}} &x\ne 0\\
    0 &x=0
\end{cases}
\]
Similar to \textbf{(a)}, with the difference of only one exponent, we observe that:
\[
\boxed{f^{(n)}(x) $ is continuous but not differentiable$}
\]

\problem{3}
Find $f'(0)$ if
\[
f(x)=\begin{cases}
    g(x)\sin \paren{ \dfrac{1}{x}} &x\ne 0\\
    0 &x=0
\end{cases}
\]
and $g(0)=g'(0)=0.$
\begin{align*}
    g'(0)
    &=\lim_{h \to 0} \frac{g(0+h)-g(0)}{h}\\
    &=\lim_{h \to 0}\frac{g(h)-0}{h}\\
    &=\lim_{h \to 0} \frac{g(h)}{h}\\
    &=0\\
    f'(0)
    &=\lim_{h \to 0}\frac{f(0+h)-f(0)}{h}\\
    &=\lim_{h \to 0} \frac{g(h)\sin \frac{1}{h}-0}{h}\\
    &=\lim_{h \to 0} \frac{g(h)}{h} \cdot \sin \frac{1}{h}\\
    &=0 \cdot \lim_{h \to 0} \sin \frac{1}{h}\\
    &=\boxed{0}
\end{align*}

\problem{4} \vspace{-0.3em}
\subproblem{a} Use the definition of differentiability to prove that function $f(x)=(x-1)^\frac{1}{5}\sin \pi x$ is differentiable at $x = 1$.
\begin{gather*}
    \lim_{h \to 0^-} \frac{f(1+h)-f(1)}{h}=\lim_{h \to 0^-} \frac{h^{\frac{1}{5}}\sin(\pi+\pi h)-0}{h}=\lim_{h \to 0^-} -\frac{\sin \pi h}{\pi h}\cdot h^\frac{1}{5}\pi=\lim_{h \to 0^-}-h^\frac{1}{5}\pi=0\\
    \lim_{h \to 0^+} \frac{f(1+h)-f(1)}{h}=\lim_{h \to 0^+} \frac{h^{\frac{1}{5}}\sin(\pi+\pi h)-0}{h}=\lim_{h \to 0^+} -\frac{\sin \pi h}{\pi h}\cdot h^\frac{1}{5}\pi=\lim_{h \to 0^+}-h^\frac{1}{5}\pi=0\\
    f'(1^-)=f'(1^+)=0 \implies f'(1)=0\\
    \boxed{f(x) $ is differentiable at $ x=1}
\end{gather*}

\subproblem{b} Use the definition of continuity to prove that $f'(x)$ is continuous at $x = 1$.
\begin{gather*}
    f'(x)= \frac{1}{5(x-1)^{\frac{4}{5}}} \sin \pi x+(x-1)^\frac{1}{5} \pi \cos \pi x\\
    \limit{1^-} f'(x)= 
    \limit{1^-} \frac{1}{5(0^-)^{\frac{4}{5}}} \sin \pi+(0^-)^\frac{1}{5} \pi \cos \pi=
    \limit{1^-}0+0=0\\
    \limit{1^+} f'(x)= 
    \limit{1^+} \frac{1}{5(0^+)^{\frac{4}{5}}} \sin \pi+(0^+)^\frac{1}{5} \pi \cos \pi=
    \limit{1^+}0+0=0\\
    f'(1^-)=f'(1^+)=0 \implies f'(1)=0\\
    \boxed{f'(x) $ is continuous at $ x=1}
\end{gather*}

\problem{5}
Find the intervals where $f(x)=x^2|x|$ is thrice differentiable.
\begin{align*}
    f(x)&=\begin{cases}
        x^3 &x \ge 0\\
        -x^3 &x<0
    \end{cases}\\
    f'(x)&=\begin{cases}
        3x^2 &x > 0\\
        0 &x=0\\
        -3x^2 &x<0
    \end{cases}\\
    f''(x)&=\begin{cases}
        6x &x > 0\\
        0 &x=0\\
        -6x &x<0
    \end{cases}\\
    f'''(x)&=\begin{cases}
        6 &x > 0\\
        -6 &x<0
    \end{cases}
\end{align*}
\[
\boxed{x\ne 0}
\]

\end{document}

